\documentclass[11pt,a4paper]{article}

\usepackage[T1]{fontenc}
\usepackage[utf8]{inputenc}
\usepackage[brazil]{babel}
\usepackage{lmodern}
\usepackage{microtype}
\usepackage[margin=2cm]{geometry}
\usepackage{hyperref}
\usepackage{titlesec}
\usepackage{xcolor}

\hypersetup{
  colorlinks=true,
  linkcolor=black,
  urlcolor=black,
  citecolor=black,
  pdfborder={0 0 0}
}

\definecolor{sectioncolor}{HTML}{111111}
\definecolor{accentcolor}{HTML}{1A1A1A}

\pagestyle{empty}
\setlength{\parindent}{0pt}
\setlength{\parskip}{5pt}
\linespread{1.05}
\selectfont

\titleformat{\section}
  {\bfseries\large\color{sectioncolor}}
  {}{0pt}{}[\vspace{-2pt}{\color{sectioncolor}\titlerule[0.8pt]}]
\titlespacing*{\section}{0pt}{12pt}{8pt}

\newcommand{\cvrole}[4]{%
  \vspace{4pt}%
  {\bfseries\color{accentcolor}#1}\hfill{\itshape #2}\\[-2pt]%
  {\itshape #3}\hfill #4\\[6pt]%
}

\begin{document}

\begin{center}
  {\fontsize{20}{24}\selectfont\bfseries Kamille Hartmann Maccagnan}\\[8pt]
  {\large Sucesso do Cliente \textbar{} Gestão de Carteira \textbar{} Análise de Dados}\\[10pt]
  Campo Comprido, Curitiba-PR \enspace\textbullet\enspace (41) 9 9928-4424 \enspace\textbullet\enspace
  \href{mailto:kamillehartmannm@gmail.com}{kamillehartmannm@gmail.com}\\[3pt]
  \href{https://linkedin.com/in/kamille-hartmann-maccagnan/}{linkedin.com/in/kamille-hartmann-maccagnan/}
\end{center}

\vspace{6pt}

\section{Resumo Profissional}

Profissional com experiência em suporte consultivo, gestão de demandas e monitoramento de indicadores em ambientes de varejo, foodservice e operações multicanal. Atua na condução de treinamentos, padronização de processos e apoio operacional a redes de lojas, com foco em adoção de soluções, melhoria de engajamento e redução de falhas no dia a dia do cliente. Vivência prática no setor de restaurantes, com análise de performance operacional, resultados e indicadores de negócio. Domínio de Excel, SQL, Power BI e Google Workspace, com perfil comunicativo, empático, organizado e orientado à resolução de problemas com autonomia.

\section{Experiência Profissional}

\cvrole{Anjuss}{Analista de Suporte de TI}{07/2025 -- Atual}

Atuo no suporte consultivo a usuários e na operação da rede de lojas, com organização de demandas, orientações funcionais e acompanhamento de casos que impactam a rotina dos clientes internos. Lidero treinamentos e capacitação de novos funcionários, elaboro POPs e documentações operacionais para lojas e equipe de T.I., promovendo adoção correta dos sistemas e padronização de processos. Desenvolvo dashboards em Power BI para monitoramento de indicadores operacionais e desempenho, priorizando riscos e oportunidades de melhoria com base em dados. Atuo na interface entre usuários e áreas técnicas, conduzindo planos de ação e automações com Power Automate para aumentar eficiência, consistência na experiência e qualidade das entregas.

\cvrole{HORSE do Brasil}{Estagiária de TI (IS/IT LATAM)}{07/2024 -- 07/2025}

Atuei no acompanhamento de demandas e projetos em ambiente multinacional, com participação em frentes de PMO, suporte e governança de dados. Utilizei o ServiceNow para gestão de chamados e acompanhamento de casos, apoiando a organização de reincidências e o alinhamento entre usuários e áreas técnicas. Elaborei relatórios executivos e dashboards para monitoramento de performance, atuando como ponte entre negócio, fornecedores e stakeholders internacionais em reuniões multilíngues. Participei da documentação de processos, análise de riscos e melhoria contínua, com postura consultiva na tradução de necessidades e acompanhamento de entregas.

\cvrole{Restaurante Familiar}{Analista de Dados e Suporte Técnico}{01/2023 -- 07/2024}

Atuei no setor de foodservice com análise de vendas, performance operacional e resultados financeiros para apoio à gestão do negócio. Estruturei do zero a base de dados em Excel e dashboards em Power BI para acompanhamento de indicadores de desempenho, identificando oportunidades de melhoria na operação do restaurante. Realizei suporte técnico e alinhamento entre áreas operacionais e administrativas, com foco em eficiência de processos, confiabilidade de informações e tomada de decisão orientada a dados. Experiência direta com desafios da operação gastronômica, incluindo rotina de atendimento, controle de resultados e melhoria contínua.

\section{Competências Técnicas}

\textbf{Sucesso do Cliente e Relacionamento:} Suporte consultivo, atendimento ao cliente, onboarding e treinamento de usuários, gestão de demandas, comunicação orientada ao cliente, escuta ativa, mapeamento de necessidades, gestão de relacionamento, padronização de processos, elaboração de POPs e playbooks operacionais.

\textbf{Análise de Dados e Indicadores:} Monitoramento de KPIs, análise de dados, relatórios gerenciais, dashboards em Power BI, Excel, SQL, Google Sheets, identificação de riscos e oportunidades com base em indicadores, acompanhamento de performance operacional.

\textbf{Integração e Resolução de Problemas:} Interface entre suporte, operações e áreas técnicas, planos de ação, análise de causas operacionais, melhoria contínua, autonomia na condução de casos de média complexidade.

\textbf{Ferramentas:} Power BI, ServiceNow, Google Workspace, Power Automate, Trello, Monday, Pacote Office.

\textbf{Metodologias:} Scrum, Kanban.

\textbf{Idiomas:} Inglês Avançado \enspace\textbullet\enspace Espanhol Básico \enspace\textbullet\enspace Coreano Básico

\section{Projetos e Conquistas}

1º lugar no Transformation Day Renault 2023 com solução orientada a inovação, dados e melhoria de processos. Participação no Hackathon Bosch 2023 com foco em resolução de problemas de negócio. Desenvolvimento de dashboards para monitoramento de indicadores e suporte à tomada de decisão. Estruturação de KPIs para análise de performance operacional. Automação de fluxos de dados e relatórios com Power Automate integrado ao Power BI.

\section{Formação Acadêmica}

\textbf{PUCPR} \hfill Curitiba, PR\\
Graduação em Gestão da Tecnologia da Informação \hfill Previsão de conclusão: 06/2026\\[6pt]

\textbf{Escola Conquer} \hfill Curitiba, PR\\
Pós-graduação em Liderança e Gestão em Tecnologia \hfill Conclusão: 2024\\[6pt]

\textbf{Universidade Tuiuti do Paraná} \hfill Curitiba, PR\\
Graduação em Medicina Veterinária \hfill 06/2019\\[6pt]

\textbf{Qualittas} \hfill Curitiba, PR\\
Pós-graduação em Clínica Médica e Cirúrgica de Animais Exóticos e Pets Não Convencionais \hfill 12/2022

\end{document}
